\section{Introduction}
Bus operators need to know where time is spent along a route. A trip total supports comparisons of operating efficiency, but cannot show whether a particular delay occurred near the beginning or the end. Interval running times support bottleneck screening and comparisons across departure periods. Obtaining complete local records, however, requires continuous position observations and reliable stop matching. Sampling intervals and missing records leave some parts of a trip unobserved. Estimating segment times from a known total and limited local observations is therefore a practical need in travel-time inference\cite{hellinga2008,hofleitner2011,yin2015entryexit}.

We study segment running times for completed trips. A segment denotes a directed inter-stop interval, and the known total is the sum of interval running times, excluding dwell. This differs from the elapsed time between endpoint timestamps, which can include dwell and other non-running time. Using such a measurement requires separating those components and aligning the interval boundaries. Our experiments sum public running-time labels to evaluate interval estimation under an exact known total.

Consider a trip through three consecutive intervals with a total running time of ten minutes. Both $(2,3,5)$ and $(5,3,2)$ minutes remain feasible. A training label of three minutes at the middle position still cannot distinguish these allocations; current-trip local labels are unavailable at inference. The two allocations identify different intervals as the most time-consuming and can lead to different operating decisions. Records of the same intervals in other trips help distinguish the alternatives. Yet changes in departure period or route context can make a historical average unrepresentative of the current trip. This example calls for both information sharing across repeated intervals and an estimate of how the current trip departs from that shared pattern.

Earlier methods use free-flow information, local time distributions, or statistical traffic-flow models to inform total-time decomposition\cite{hellinga2008,hofleitner2011}. When paths are incomplete, link inference must also recover the traveled path\cite{yin2015entryexit,sun2022prltte}. We focus on known ordered paths with recurring interval identities and a subset of observed singleton labels. Path recovery is unnecessary in this setting. The remaining question is how to combine information shared across trips with current-trip variation while making interval outputs consistent with the known total.

A segment mean shares repeated observations but treats context-varying durations as fixed. Our direct sequence baseline also shares interval embeddings and uses current-trip context, but outputs positive scores without a separate redistribution branch. We separate these roles: a shared branch estimates contextual base times, and a sequence branch proposes bounded redistribution before proportional rescaling to the known total. The contribution is this allocation parameterization and its analysis under sparse labels; an accuracy advantage over direct prediction is an empirical question.

The experiments further delimit the role of this design. Segment means remain competitive when local labels are scarce. At larger budgets, our method improves on segment means but shows no consistent advantage over same-input direct prediction. Disabling trip correction increases the error of the same model, indicating that the correction helps the base estimate; this comparison does not establish that the complete method is better than direct prediction. The same number of labels also produces substantially different errors when it covers fewer interval identities. In a control that preserves the identity set and the count for each identity, randomizing label positions does not recover the accuracy obtained with broader coverage. A known total rules out allocations that violate the sum, but cannot by itself determine which local allocation is correct; useful local estimates also depend on how well repeated intervals are represented in the supervision.
